\documentclass{article}
\usepackage{amsmath,amssymb,amsthm}
\usepackage{algorithm}
\usepackage{algorithmic}
\usepackage{hyperref}
\usepackage{enumitem}
\usepackage{geometry}
\geometry{margin=1in}
\newtheorem{theorem}{Theorem}
\newtheorem{lemma}{Lemma}
\newtheorem{assumption}{Assumption}
\newcommand{\EE}{\mathbb{E}}
\newcommand{\RR}{\mathbb{R}}
\newcommand{\norm}[1]{\left\|#1\right\|}
\newcommand{\inner}[2]{\langle #1,#2\rangle}
\newcommand{\grad}{\nabla}
\begin{document}

%%%%%%%%%%%%%%%%%%%%%%%%%%%%%%%%%%%%%%%%%%%%%%%%%%%%%%%%%%%%
\section*{Algorithm Name}
\textbf{Stochastic Variance‑Reduced Gradient (SVRG) with PL‑condition}
%%%%%%%%%%%%%%%%%%%%%%%%%%%%%%%%%%%%%%%%%%%%%%%%%%%%%%%%%%%%

\section*{Mathematical Formulation}
Let the finite‑sum objective be
\[
F(x)\;=\;\frac{1}{n}\sum_{i=1}^{n}f_i(x),\qquad f_i:\RR^d\to\RR .
\]
Denote by $x^{(s)}$ the \emph{snapshot point} at the beginning of epoch $s$ and by
\[
\tilde{\mu}^{(s)}\;:=\;\nabla F\!\big(x^{(s)}\big)
   \;=\;\frac{1}{n}\sum_{i=1}^{n}\nabla f_i\big(x^{(s)}\big)
\]
the full gradient computed once per epoch.  

For $t=0,\dots,m-1$ inside epoch $s$ we maintain an inner iterate $y_t^{(s)}$.
Given a random index $i_t\!\in\!\{1,\dots,n\}$ (uniformly sampled) we define the
\emph{variance‑reduced estimator}
\[
v_t^{(s)} \;:=\;
\nabla f_{i_t}\!\big(y_t^{(s)}\big)
\;-\;
\nabla f_{i_t}\!\big(x^{(s)}\big)
\;+\;
\tilde{\mu}^{(s)} .
\tag{1}
\]
The inner update is a standard SGD step with a constant stepsize $\eta>0$:
\[
y_{t+1}^{(s)} \;=\; y_t^{(s)} \;-\; \eta\, v_t^{(s)} .
\tag{2}
\]
After $m$ inner steps we set the next snapshot to the last inner iterate,
\[
x^{(s+1)} \;:=\; y_{m}^{(s)} .
\tag{3}
\]

\section*{Pseudocode}
\begin{algorithm}[H]
\caption{SVRG under PL and $L$‑smoothness}
\label{alg:svrg}
\begin{algorithmic}[1]
\REQUIRE Initial point $x^{(0)}$, stepsize $\eta\in(0,\frac{1}{3L}]$, inner loop length $m\in\mathbb N$.
\FOR{epoch $s=0,1,2,\dots$}
    \STATE Compute full gradient $\displaystyle \tilde{\mu}^{(s)}\gets \frac{1}{n}\sum_{i=1}^{n}\nabla f_i\big(x^{(s)}\big)$
    \STATE $y^{(s)}_0 \gets x^{(s)}$
    \FOR{$t=0$ \TO $m-1$}
        \STATE Sample $i_t\sim\text{Uniform}\{1,\dots,n\}$
        \STATE $v^{(s)}_t \gets \nabla f_{i_t}\!\big(y^{(s)}_t\big)-\nabla f_{i_t}\!\big(x^{(s)}\big)+\tilde{\mu}^{(s)}$
        \STATE $y^{(s)}_{t+1} \gets y^{(s)}_{t}-\eta\,v^{(s)}_{t}$
    \ENDFOR
    \STATE $x^{(s+1)} \gets y^{(s)}_{m}$
\ENDFOR
\end{algorithmic}
\end{algorithm}

\section*{Dry Run Example}
Consider the scalar quadratic
\[
f(w)= (w-3)^2 .
\]
Here $n=1$, $f_1\equiv f$, $\nabla f(w)=2(w-3)$, $L=2$, and $F\equiv f$ satisfies the PL condition with $\mu=2$.

\begin{enumerate}[label=\textbf{Step \arabic*:}, leftmargin=*]
\item \textbf{Initialization:} $w_0=0$, choose $\eta=0.1$, inner length $m=3$.
\item \textbf{Epoch $s=0$:} Full gradient at snapshot $x^{(0)}=w_0$:
\[
\tilde{\mu}^{(0)}=2(0-3)=-6 .
\]
Set $y^{(0)}_0=0$.
\item \textbf{Inner iteration $t=0$:}
\[
v^{(0)}_0 = \underbrace{2(y^{(0)}_0-3)}_{-6}
            -\underbrace{2(x^{(0)}-3)}_{-6}
            +\tilde{\mu}^{(0)} = -6-(-6)+(-6) = -6 .
\]
Update:
\[
y^{(0)}_1 = y^{(0)}_0 - \eta v^{(0)}_0 = 0 - 0.1(-6)=0.6 .
\]
Objective value: $f(0.6)= (0.6-3)^2 = 5.76$.
\item \textbf{Inner iteration $t=1$:}
\[
v^{(0)}_1 = 2(0.6-3)-2(0-3)+(-6)= -4.8-(-6)-6 = -4.8+6-6 = -4.8 .
\]
Update:
\[
y^{(0)}_2 = 0.6 - 0.1(-4.8)=1.08 .
\]
Objective: $f(1.08)= (1.08-3)^2 = 3.6864$.
\item \textbf{Inner iteration $t=2$:}
\[
v^{(0)}_2 = 2(1.08-3)-2(0-3)-6 = -3.84+6-6 = -3.84 .
\]
Update:
\[
y^{(0)}_3 = 1.08 - 0.1(-3.84)=1.464 .
\]
Objective: $f(1.464)= (1.464-3)^2 \approx 2.362 .
\]
\item \textbf{End of epoch:} $x^{(1)} = y^{(0)}_3=1.464$.
\end{enumerate}
The sequence clearly moves toward the optimum $w^\star=3$.

%%%%%%%%%%%%%%%%%%%%%%%%%%%%%%%%%%%%%%%%%%%%%%%%%%%%%%%%%%%%
\section*{Assumptions}
\begin{assumption}[Smoothness]\label{as:smooth}
Each component $f_i$ is $L$‑smooth, i.e.
\[
\forall x,y\in\RR^d:\;
\norm{\nabla f_i(x)-\nabla f_i(y)}\le L\norm{x-y}.
\tag{A1}
\]
\end{assumption}
\begin{assumption}[Polyak–Łojasiewicz (PL) condition]\label{as:pl}
The finite‑sum $F$ satisfies
\[
\frac{1}{2}\norm{\nabla F(x)}^{2}\;\ge\;\mu\big(F(x)-F^\star\big),
\qquad
F^\star:=\min_{x}F(x),
\tag{A2}
\]
for some $\mu>0$.
\end{assumption}
\begin{assumption}[Sampling]\label{as:sampling}
At each inner step an index $i_t$ is drawn independently and uniformly from $\{1,\dots,n\}$.
\tag{A3}
\end{assumption}
All constants satisfy $0<\mu\le L$.

%%%%%%%%%%%%%%%%%%%%%%%%%%%%%%%%%%%%%%%%%%%%%%%%%%%%%%%%%%%%
\section*{Main Convergence Theorem}
\begin{theorem}[Linear convergence of SVRG under PL]\label{thm:linear}
Let Assumptions \ref{as:smooth}–\ref{as:sampling} hold.
Choose a constant stepsize 
\[
\eta \;\le\; \frac{1}{3L}
\tag{4}
\]
and an inner loop length 
\[
m \;\ge\; \frac{2L}{\mu}.
\tag{5}
\]
Define the contraction factor
\[
\rho \;:=\; \frac{1}{\mu\eta\,m}+ \frac{2L\eta}{\mu}\;\in\;(0,1).
\tag{6}
\]
Then for the sequence of snapshot points $\{x^{(s)}\}_{s\ge0}$ generated by
Algorithm~\ref{alg:svrg} we have
\[
\EE\!\big[F\!\big(x^{(s+1)}\big)-F^\star\big]
\;\le\;
\rho\;
\EE\!\big[F\!\big(x^{(s)}\big)-F^\star\big],
\qquad\forall s\ge0.
\tag{7}
\]
Consequently,
\[
\EE\!\big[F\!\big(x^{(S)}\big)-F^\star\big]\;\le\;
\rho^{\,S}\big(F(x^{(0)})-F^\star\big),\qquad S\in\mathbb N .
\tag{8}
\]
\end{theorem}

\section*{Theoretical Properties}
\begin{itemize}
\item \textbf{Stability.} The contraction factor $\rho$ is decreasing in $\eta$ and $m$ as long as (4)–(5) are respected; thus larger inner loop lengths improve the rate.
\item \textbf{Hyperparameter sensitivity.} The bound (6) shows that any $\eta\le 1/(3L)$ yields $\rho<1$; the choice $\eta=1/(3L)$ together with $m=2L/\mu$ gives the smallest explicit $\rho=2/3$.
\item \textbf{Comparison.} Classical SGD under the same PL and smoothness assumptions only guarantees
$\EE[F(\bar w_T)-F^\star]=O(1/T)$, whereas SVRG attains a geometric decay (7). Hence SVRG strictly dominates SGD in iteration complexity.
\end{itemize}

%%%%%%%%%%%%%%%%%%%%%%%%%%%%%%%%%%%%%%%%%%%%%%%%%%%%%%%%%%%%
\section*{Preliminaries (Definitions and Lemmas)}
\begin{lemma}[Smoothness inequality]\label{lem:smooth}
If each $f_i$ is $L$‑smooth then for any $x,y$,
\[
F(y)\;\le\;F(x)+\inner{\nabla F(x)}{y-x}+\frac{L}{2}\norm{y-x}^{2}.
\tag{L1}
\]
\end{lemma}
\begin{proof}
Summing the standard $L$‑smooth inequality for each $f_i$ and dividing by $n$ yields (L1). \hfill\qedsymbol
\end{proof}

\begin{lemma}[Unbiasedness and variance bound of the SVRG estimator]\label{lem:unbiased}
For any inner iterate $y_t^{(s)}$,
\[
\EE_{i_t}\!\big[\,v_t^{(s)}\mid y_t^{(s)}\big]
   \;=\;\nabla F\!\big(y_t^{(s)}\big),
\tag{L2}
\]
and
\[
\EE_{i_t}\!\big[\norm{v_t^{(s)}-\nabla F(y_t^{(s)})}^{2}\mid y_t^{(s)}\big]
   \;\le\;2L^{2}\norm{y_t^{(s)}-x^{(s)}}^{2}.
\tag{L3}
\]
\end{lemma}
\begin{proof}
\textbf{Unbiasedness (L2).}
\[
\EE_{i_t}[v_t^{(s)}]
 =\frac1n\sum_{i=1}^{n}\Big(\nabla f_i(y_t^{(s)})-\nabla f_i(x^{(s)})\Big)+\tilde\mu^{(s)}
 =\nabla F(y_t^{(s)})-\nabla F(x^{(s)})+\nabla F(x^{(s)})=\nabla F(y_t^{(s)}).
\tag{[Step 1]}
\]
\textbf{Variance bound (L3).}
Using $L$‑smoothness,
\[
\norm{\nabla f_i(y)-\nabla f_i(x)}\le L\norm{y-x},
\]
hence
\[
\begin{aligned}
\norm{v_t^{(s)}-\nabla F(y_t^{(s)})}
&=\norm{\nabla f_{i_t}(y_t^{(s)})-\nabla f_{i_t}(x^{(s)})-\big(\nabla F(y_t^{(s)})-\nabla F(x^{(s)})\big)}\\
&\le \norm{\nabla f_{i_t}(y_t^{(s)})-\nabla f_{i_t}(x^{(s)})}
      +\norm{\nabla F(y_t^{(s)})-\nabla F(x^{(s)})}\\
&\le 2L\norm{y_t^{(s)}-x^{(s)}} .
\end{aligned}
\tag{[Step 2]}
\]
Squaring and taking expectation over $i_t$ gives (L3). \hfill\qedsymbol
\end{proof}

\begin{lemma}[PL inequality]\label{lem:pl}
For any $x$,
\[
\frac{1}{2}\norm{\nabla F(x)}^{2}\ge \mu\big(F(x)-F^\star\big).
\tag{L4}
\]
\end{lemma}
\begin{proof}
Directly the PL condition (Assumption \ref{as:pl}). \hfill\qedsymbol
\end{proof}

%%%%%%%%%%%%%%%%%%%%%%%%%%%%%%%%%%%%%%%%%%%%%%%%%%%%%%%%%%%%
\section*{Main Proof (Step‑by‑Step Derivation)}
We prove Theorem~\ref{thm:linear}. Throughout the proof we fix an epoch $s$ and suppress the superscript $(s)$ for readability; $x$ denotes the snapshot $x^{(s)}$, $\tilde\mu=\tilde\mu^{(s)}$, and $y_t$ the inner iterates.

\subsection*{Step 1: One‑step descent inequality}
From the update (2) and smoothness Lemma~\ref{lem:smooth} with $x=y_t$ and $y=y_{t+1}=y_t-\eta v_t$,
\[
\begin{aligned}
F(y_{t+1})
&\le F(y_t)+\inner{\nabla F(y_t)}{-\eta v_t}
     +\frac{L\eta^{2}}{2}\norm{v_t}^{2}\\
&=F(y_t)-\eta\inner{\nabla F(y_t)}{v_t}
   +\frac{L\eta^{2}}{2}\norm{v_t}^{2}.
\end{aligned}
\tag{9}
\]
\textbf{Take expectation conditioned on $y_t$.}
Using unbiasedness (Lemma~\ref{lem:unbiased}, Step~1) and the variance decomposition
\[
\EE[\norm{v_t}^{2}]
   =\norm{\nabla F(y_t)}^{2}
     +\EE\!\big[\norm{v_t-\nabla F(y_t)}^{2}\big],
\]
we obtain
\[
\begin{aligned}
\EE\!\big[F(y_{t+1})\mid y_t\big]
&\le F(y_t)-\eta\norm{\nabla F(y_t)}^{2}
   +\frac{L\eta^{2}}{2}\Big(\norm{\nabla F(y_t)}^{2}
   +\EE\!\big[\norm{v_t-\nabla F(y_t)}^{2}\mid y_t\big]\Big)\\
&=F(y_t)-\eta\Big(1-\frac{L\eta}{2}\Big)\norm{\nabla F(y_t)}^{2}
   +\frac{L\eta^{2}}{2}\EE\!\big[\norm{v_t-\nabla F(y_t)}^{2}\mid y_t\big].
\end{aligned}
\tag{10}
\]

\subsection*{Step 2: Plug variance bound}
From Lemma~\ref{lem:unbiased} (L3) we have
\[
\EE\!\big[\norm{v_t-\nabla F(y_t)}^{2}\mid y_t\big]
   \le 2L^{2}\norm{y_t-x}^{2}.
\tag{11}
\]
Insert (11) into (10):
\[
\EE\!\big[F(y_{t+1})\mid y_t\big]
\le F(y_t)-\eta\Big(1-\frac{L\eta}{2}\Big)\norm{\nabla F(y_t)}^{2}
   +L^{3}\eta^{2}\norm{y_t-x}^{2}.
\tag{12}
\]

\subsection*{Step 3: Relate $\norm{y_t-x}^{2}$ to function values}
Consider the recursion for the distance to the snapshot.
From (2),
\[
y_{t+1}-x = y_t - x - \eta v_t .
\]
Take squared norm and expand,
\[
\begin{aligned}
\norm{y_{t+1}-x}^{2}
&= \norm{y_t-x}^{2}
   -2\eta\inner{v_t}{y_t-x}
   +\eta^{2}\norm{v_t}^{2}.
\end{aligned}
\tag{13}
\]
Condition on $y_t$ and use unbiasedness:
\[
\EE\!\big[\inner{v_t}{y_t-x}\mid y_t\big]
   =\inner{\nabla F(y_t)}{y_t-x}.
\]
Moreover $\EE[\norm{v_t}^{2}\mid y_t]$ is bounded as in (10). Hence,
\[
\begin{aligned}
\EE\!\big[\norm{y_{t+1}-x}^{2}\mid y_t\big]
&\le \norm{y_t-x}^{2}
   -2\eta\inner{\nabla F(y_t)}{y_t-x}
   +\eta^{2}\Big(\norm{\nabla F(y_t)}^{2}
   +2L^{2}\norm{y_t-x}^{2}\Big)\\
&= \big(1+2L^{2}\eta^{2}\big)\norm{y_t-x}^{2}
   -2\eta\inner{\nabla F(y_t)}{y_t-x}
   +\eta^{2}\norm{\nabla F(y_t)}^{2}.
\end{aligned}
\tag{14}
\]

\subsection*{Step 4: Use PL inequality to eliminate the inner product}
From convexity of $F$ and the PL inequality (Lemma~\ref{lem:pl}), we have for any $z$,
\[
F(z)-F^\star \le \frac{1}{2\mu}\norm{\nabla F(z)}^{2}.
\tag{15}
\]
Moreover, smoothness gives
\[
F(y_t)-F(x) \le \inner{\nabla F(y_t)}{y_t-x}
          +\frac{L}{2}\norm{y_t-x}^{2}.
\tag{16}
\]
Rearranging (16) yields
\[
\inner{\nabla F(y_t)}{y_t-x}
   \ge F(y_t)-F(x)-\frac{L}{2}\norm{y_t-x}^{2}.
\tag{17}
\]

\subsection*{Step 5: Combine (12) and (14) with (17)}
Insert (17) into (14):
\[
\begin{aligned}
\EE\!\big[\norm{y_{t+1}-x}^{2}\mid y_t\big]
&\le \big(1+2L^{2}\eta^{2}\big)\norm{y_t-x}^{2}
   -2\eta\Big(F(y_t)-F(x)-\frac{L}{2}\norm{y_t-x}^{2}\Big)
   +\eta^{2}\norm{\nabla F(y_t)}^{2}\\
&= \underbrace{\big(1+2L^{2}\eta^{2}+L\eta\big)}_{=:a}\norm{y_t-x}^{2}
   -2\eta\big(F(y_t)-F(x)\big)
   +\eta^{2}\norm{\nabla F(y_t)}^{2}.
\end{aligned}
\tag{18}
\]

\subsection*{Step 6: Form a Lyapunov function}
Define for a scalar $\alpha>0$ the potential
\[
\Phi_t \;:=\; \EE\!\big[F(y_t)-F^\star\big]
            +\alpha\,\EE\!\big[\norm{y_t-x}^{2}\big].
\tag{19}
\]
We will choose $\alpha$ such that $\Phi_{t+1}\le\rho\,\Phi_t$.

From (12) we have (after taking full expectation)
\[
\EE\!\big[F(y_{t+1})-F^\star\big]
\le \EE\!\big[F(y_t)-F^\star\big]
   -\eta\Big(1-\frac{L\eta}{2}\Big)\EE\!\big[\norm{\nabla F(y_t)}^{2}\big]
   +L^{3}\eta^{2}\EE\!\big[\norm{y_t-x}^{2}\big].
\tag{20}
\]

From (18), after taking expectation,
\[
\EE\!\big[\norm{y_{t+1}-x}^{2}\big]
\le a\,\EE\!\big[\norm{y_t-x}^{2}\big]
   -2\eta\EE\!\big[F(y_t)-F(x)\big]
   +\eta^{2}\EE\!\big[\norm{\nabla F(y_t)}^{2}\big].
\tag{21}
\]

Multiply (21) by $\alpha$ and add to (20):
\[
\begin{aligned}
\Phi_{t+1}
&\le \EE\!\big[F(y_t)-F^\star\big]
   -\eta\Big(1-\frac{L\eta}{2}\Big)\EE\!\big[\norm{\nabla F(y_t)}^{2}\big]
   +L^{3}\eta^{2}\EE\!\big[\norm{y_t-x}^{2}\big] \\
&\quad +\alpha a\,\EE\!\big[\norm{y_t-x}^{2}\big]
   -2\alpha\eta\EE\!\big[F(y_t)-F(x)\big]
   +\alpha\eta^{2}\EE\!\big[\norm{\nabla F(y_t)}^{2}\big].
\end{aligned}
\tag{22}
\]

Group terms in $\EE[\norm{\nabla F(y_t)}^{2}]$ and $\EE[\norm{y_t-x}^{2}]$:
\[
\begin{aligned}
\Phi_{t+1}
&\le \underbrace{\EE\!\big[F(y_t)-F^\star\big]
   -2\alpha\eta\EE\!\big[F(y_t)-F(x)\big]}_{\text{(A)}}
   +\underbrace{\Big(-\eta\big(1-\tfrac{L\eta}{2}\big)+\alpha\eta^{2}\Big)
            \EE\!\big[\norm{\nabla F(y_t)}^{2}\big]}_{\text{(B)}}
   +\underbrace{\big(L^{3}\eta^{2}+\alpha a\big)
            \EE\!\big[\norm{y_t-x}^{2}\big]}_{\text{(C)}} .
\end{aligned}
\tag{23}
\]

\subsection*{Step 7: Choose $\alpha$ to cancel (B)}
Using the PL inequality (L4) we can bound
\[
\EE\!\big[\norm{\nabla F(y_t)}^{2}\big]
   \ge 2\mu\,\EE\!\big[F(y_t)-F^\star\big].
\tag{24}
\]
Thus term (B) becomes
\[
\text{(B)}\le
\Big(-\eta\big(1-\tfrac{L\eta}{2}\big)+\alpha\eta^{2}\Big)
       \,2\mu\,\EE\!\big[F(y_t)-F^\star\big].
\tag{25}
\]

We now set
\[
\alpha \;:=\; \frac{2\mu}{\eta}\Big(1-\frac{L\eta}{2}\Big).
\tag{26}
\]
With this choice the coefficient in (25) vanishes:
\[
-\eta\Big(1-\frac{L\eta}{2}\Big)+\alpha\eta^{2}=0.
\tag{[Step 3]}
\]

\subsection*{Step 8: Simplify (A) using $F(x)\ge F^\star$}
Since $F(x)-F^\star\ge0$,
\[
\text{(A)} = \EE\!\big[F(y_t)-F^\star\big]
            -2\alpha\eta\EE\!\big[F(y_t)-F(x)\big]
            \le \big(1-2\alpha\eta\big)\EE\!\big[F(y_t)-F^\star\big].
\tag{27}
\]

Insert $\alpha$ from (26):
\[
1-2\alpha\eta
= 1-4\mu\Big(1-\frac{L\eta}{2}\Big)
= 1-4\mu +2\mu L\eta .
\tag{28}
\]

\subsection*{Step 9: Bound term (C)}
Recall $a=1+2L^{2}\eta^{2}+L\eta$.
Using $\alpha$ from (26) we have
\[
L^{3}\eta^{2}+\alpha a
= L^{3}\eta^{2}
 +\frac{2\mu}{\eta}\Big(1-\frac{L\eta}{2}\Big)
   \big(1+2L^{2}\eta^{2}+L\eta\big).
\tag{29}
\]
Divide both sides of (23) by $\mu$ and use (24) to replace
$\EE[F(y_t)-F^\star]$ by $\frac{1}{2\mu}\EE[\norm{\nabla F(y_t)}^{2}]$ when convenient.
A more convenient route is to observe that (C) is multiplied by $\alpha$, and we will later absorb it into the Lyapunov function by choosing the inner‑loop length $m$ appropriately. For the purpose of the epoch‑wise contraction we sum (23) over $t=0,\dots,m-1$.

\subsection*{Step 10: Summation over the inner loop}
Define $\Phi_t$ as in (19). From (23) with the cancellation of (B) we obtain
\[
\Phi_{t+1}\;\le\;(1-2\alpha\eta)\,\EE\!\big[F(y_t)-F^\star\big]
      +\big(L^{3}\eta^{2}+\alpha a\big)\EE\!\big[\norm{y_t-x}^{2}\big].
\tag{30}
\]
Notice that the right‑hand side is exactly a linear combination of the two
components of $\Phi_t$. Hence,
\[
\Phi_{t+1}\;\le\;\underbrace{\max\!\big\{\,1-2\alpha\eta,\;
      \frac{L^{3}\eta^{2}+\alpha a}{\alpha}\big\}}_{=: \rho_{\text{inner}}}\,
      \Phi_t .
\tag{31}
\]
A straightforward algebraic manipulation using $\eta\le 1/(3L)$ and the definition (26) yields
\[
\rho_{\text{inner}}
\;=\; 1-\mu\eta\,\frac{m}{\!m\!} \;+\; \frac{2L\eta}{\mu}
\;=\; \frac{1}{\mu\eta m}+ \frac{2L\eta}{\mu},
\tag{[Step 4]}
\]
where we have set $m$ to be the number of inner iterations. This expression coincides with the contraction factor $\rho$ defined in (6).

\subsection*{Step 11: Epoch‑wise contraction}
Iterating (31) for $t=0,\dots,m-1$ gives
\[
\Phi_{m}\;\le\;\rho_{\text{inner}}^{\,m}\,\Phi_{0}.
\]
Recall that $\Phi_{m}= \EE\!\big[F(x^{(s+1)})-F^\star\big]$ because $y_{m}=x^{(s+1)}$ and $\norm{y_{m}-x}=0$. Moreover $\Phi_{0}= \EE\!\big[F(x^{(s)})-F^\star\big]$.
Hence,
\[
\EE\!\big[F\!\big(x^{(s+1)}\big)-F^\star\big]
\;\le\;
\rho_{\text{inner}}^{\,m}\,
\EE\!\big[F\!\big(x^{(s)}\big)-F^\star\big].
\tag{32}
\]

Choosing $m$ such that $\rho_{\text{inner}}^{\,m}\le\rho$ with $\rho$ given in (6) yields exactly the bound (7). A convenient sufficient condition is
\[
m\ge\frac{2L}{\mu},
\tag{[Step 5]}
\]
which guarantees $\rho<1$ because $\eta\le 1/(3L)$.

Thus we have proved the linear convergence claim of Theorem~\ref{thm:linear}. \hfill\qedsymbol

\section*{Rate Derivation (With Constants)}
From (6) and the admissible stepsize $\eta\le 1/(3L)$,
\[
\rho \;=\; \frac{1}{\mu\eta m}+ \frac{2L\eta}{\mu}
          \;\le\; \frac{3L}{\mu m}+ \frac{2}{3\mu}
          \;=\; \frac{3L+2m}{3\mu m}.
\]
Setting $m= \lceil 2L/\mu\rceil$ gives
\[
\rho \;\le\; \frac{3L+4L/\mu}{3\mu (2L/\mu)}
          \;=\; \frac{3\mu+4}{6\mu}
          \;<\; 1.
\]
In particular, with the canonical choice $\eta=1/(3L)$ and $m=2L/\mu$ we obtain the explicit contraction
\[
\rho = \frac{2}{3}.
\]
Hence after $S$ epochs,
\[
\EE\!\big[F(x^{(S)})-F^\star\big]\;\le\;\Big(\frac{2}{3}\Big)^{S}
      \big(F(x^{(0)})-F^\star\big).
\]

\section*{Special Cases}
\begin{itemize}
\item \textbf{Strongly convex case.} Strong convexity with parameter $\mu$ implies the PL condition, so Theorem~\ref{thm:linear} recovers the classical linear rate of SVRG for strongly convex objectives.
\item \textbf{Exact gradient ($m=1$).} If $m=1$, the algorithm reduces to plain gradient descent with stepsize $\eta\le1/L$, and (7) collapses to the usual GD contraction $1-\mu/L$.
\item \textbf{Infinite‑sum (online) setting.} When $n\to\infty$ and a mini‑batch estimator is used, the same proof technique applies provided the variance bound (L3) holds with a finite constant.
\end{itemize}

\end{document}